\section{A connection from successive time averages}
\label{sec:connection}

Proposition~\ref{cy:trace-identity} allows us to compute the trace with any
connection satisfying Definition~\ref{cy:connection}.  We now construct a
connection whose curvature can be estimated at several time scales.  The
construction first gives a one-form $B$ with uniformly positive contraction
$B(V)$; division by this contraction will produce the connection.  The
summary channels of Proposition~\ref{ch:summary} make the contributions of
successive scales telescope.  The final contribution uses the twist by $g$.

Throughout this section $G$ is finitely generated, $d$ is its fixed left
invariant word metric, and $g\in G$.  Constants denoted by $C$ may depend on
the word metric and on the fixed smooth functions used below.  After the
parameter choice at the end of the construction, a subscript $R$ will
permit dependence on the radius $R$, but never on the tuple, its
dimension, or the element $g$.

\subsection{Time averages and spatial summaries}

Begin with a finite number $b_*\geq2$ of averaging stages, an even ambient
dimension $n$, and positive time scales $\delta_1<\cdots<\delta_{b_*}$
satisfying
\begin{equation}
 \sum_{i=1}^j\delta_i\leq2\delta_j,
 \qquad 10\delta_{b_*}<1.
 \label{co:scale-support}
\end{equation}
The construction through positivity uses only these conditions on the
time scales.  Their precise dependence on $R$ and $n$ will be selected
afterwards, when preparing the quantitative estimates.
\emph{The same scales are kept fixed in every dimension $0\leq k\leq n$.}
Thus passing to a face never changes the averaging operation.

Choose a nonnegative, even $C^\infty$ function $\mu$, supported in
$[-1,1]$, with integral one, and put
$\mu_j(s)=\delta_j^{-1}\mu(s/\delta_j)$.  On a tuple of dimension $k$,
write its step function and its averages as
\begin{equation}
 \begin{split}
 f(u)&=h_i\quad(t_i\leq u<t_{i+1}),\qquad f(u+1)=gf(u),\\
 P_0(u)&=\delta_{f(u)},\qquad
 P_j(u)=\int_{\R}\mu_j(s)P_{j-1}(u-s)\,ds.
 \end{split}
 \label{co:averages}
\end{equation}
Zero-length intervals are omitted.  These are finitely supported
probability distributions at each time, and
$P_j(u+1)=gP_j(u)$.  If
$\eta_j=\mu_j*\cdots*\mu_1$, then
\begin{equation}
 P_j=\eta_j*P_0,\qquad
 \supp\eta_j\subset[-2\delta_j,2\delta_j],\qquad
 \|\eta_j\|_\infty\leq C/\delta_j.
 \label{co:smearing}
\end{equation}
The support bound follows from \eqref{co:scale-support}, and the supremum
bound follows by convolving $\mu_j$ with probability densities.

To summarize these averages spatially, choose $R\geq1$ and an integer
$2\leq J_*\leq b_*$.  Put
\begin{equation}
 r_j=6^{\min(j,J_*)-1}R,\qquad r=6^{J_*-1}R.
 \label{co:spatial-radii}
\end{equation}
The radius increases by a factor of six during the first $J_*-1$ steps
and is then fixed at $r$.  Use accuracies
\begin{equation}
 \epsilon_j=\epsilon_*(1+j)^{-2}\quad(j\geq1),
 \qquad 30\sum_{j\geq1}\epsilon_j\leq\frac1{10},
 \label{co:error-budget}
\end{equation}
where $0<\epsilon_*\leq10^{-3}$ is fixed sufficiently small.
The summable errors will leave a positive quantity after telescoping.

For $1\leq j<b_*$, apply Proposition~\ref{ch:summary} to $P_j(u)$ at
parameters $(r_j,\epsilon_j)$.  Use its early version when $j<J_*$ and its
later version when $j\geq J_*$.  Denote the resulting channels by $A_j,L_j$
and their prediction kernel by $\kappa_j$.  On the disjoint union of the
two label sets put
\begin{equation}
 Y_j(u)=A_j(P_j(u))\oplus L_j(P_j(u)).
 \label{co:channel-fields}
\end{equation}
Define a symmetric bilinear kernel $K_j$ to be zero between labels in the
same channel and to have value $\kappa_j(y,\omega)/2$ between the two
channels.  Thus
\begin{equation}
 K_j(Y_j(u),Y_j(u))
 =\sum_y A_j(P_j(u))_y\operatorname{pred}_j(P_j(u),y).
 \label{co:channel-pairing}
\end{equation}
The total mass of $Y_j(u)$ is at most two, and $|K_j|\leq1$.

At the first and last stages use
\begin{equation}
 \begin{array}{lll}
 Y_0=P_0,&
 K_0(h,h')=\chi_R(d(h,h')),\\[2pt]
 Y_{b_*}(u)=D(P_{b_*}(u)),&
 K_{b_*}(h,h')=\mathbf1_{h=h'},
 \end{array}
 \label{co:end-fields}
\end{equation}
where $D$ is the channel of Lemma~\ref{ch:terminal} at parameters
$(r,\epsilon_{b_*})$.  These two fields have mass at most one.  All the
fields obey $Y_j(u+1)=gY_j(u)$, and every kernel is invariant under
simultaneous left translation of its labels.

The last channel needs a multiplier to compensate for its possible
spread over many labels.  Lemma~\ref{ch:terminal} gives
$\sum_zD(P)_z^2\geq\frac14e^{-C_Dr}$ when $\phi_r(P)>0$, for a fixed
constant $C_D$.  Choose a fixed $C_*\geq C_D$ and set
\begin{equation}
 W=\exp(C_*r).
 \label{co:accuracy}
\end{equation}
Since $0\leq\phi_r\leq1-a_*=1/5$, this gives, for every accuracy
parameter in \eqref{co:error-budget},
\begin{equation}
 W\sum_zD(P)_z^2\geq\phi_r(P).
 \label{co:terminal-multiplier}
\end{equation}
We shall construct the connection under either of the assumptions
\begin{equation}
 \begin{array}{ll}
 \textnormal{(i)}&g\text{ has infinite order},\\[2pt]
 \textnormal{(ii)}&g\text{ has finite order and }
                \displaystyle\min_{h\in G}d(h,gh)>4r.
 \end{array}
 \label{co:twist-assumption}
\end{equation}
Only the final time-lag kernel and its contraction distinguish these two
cases.

\subsection{Lag pairings and their compatibility}

We next turn the fields into one-forms.  For $0\leq j<b_*$ let
$\rho_j=\mu_{j+1}*\mu_{j+1}$ and define the odd function
\begin{equation}
 U_j(s)=\operatorname{sign}(s)-2\int_0^s\rho_j(a)\,da.
 \label{co:lag-kernels}
\end{equation}
It is bounded in absolute value by one and vanishes outside
$[-2\delta_{j+1},2\delta_{j+1}]$.  In the sense of distributions,
\begin{equation}
 U_j'=2(\delta_0-\rho_j).
 \label{co:lag-derivative}
\end{equation}
Here $\delta_0$ denotes the Dirac measure at zero, rather than a time
scale.  At the final stage set
\begin{equation}
 \begin{array}{lll}
 U_{b_*}(s)=\operatorname{sign}(s),&\rho_{b_*}=0,
 &\text{in case (i)},\\[2pt]
 U_{b_*}(s)=\operatorname{sign}(s)\mathbf1_{|s|<1},&
 \rho_{b_*}=\tfrac12(\delta_{-1}+\delta_1),
 &\text{in case (ii)}.
 \end{array}
 \label{co:last-kernel}
\end{equation}
Formula \eqref{co:lag-derivative} holds at this stage as well.  Values at
the jump points have no effect on any integral below.

The differential $dY_j(u)$ in the following definition acts on the tuple
times $t_0,\ldots,t_k$, with $u$ fixed.  Write $\int_{\rm per}$ for
integration over one period, using a half-open interval when integrating
the raw impulse measure, and put
\begin{equation}
 \begin{split}
 B_j&=-\frac12\int_{\rm per}du\int_{\R}dv\,
             U_j(v-u)K_j(dY_j(u),Y_j(v)),\\
 B&=\sum_{j=0}^{b_*-1}B_j+WB_{b_*}.
 \end{split}
 \label{co:forms}
\end{equation}
For $j=0$ the differential of the step field is a measure in $u$.
Specifically, set
\begin{equation}
 J_i=\delta_{h_{i-1}}-\delta_{h_i},\qquad
 dP_0(u)=\sum_{i=0}^k\sum_{a\in\Z}
                  g^aJ_i\,\delta_{t_i+a}(u)\,dt_i.
 \label{co:impulses}
\end{equation}
The periodic convention gives $h_{-1}=g^{-1}h_k$.  Thus the coefficient of
$dt_i$ in the raw form is
\begin{equation}
 (B_0)_i=-\frac12\int_{\R}U_0(v-t_i)K_0(J_i,P_0(v))\,dv.
 \label{co:raw-coefficient}
\end{equation}
The expression in \eqref{co:impulses} is a distributional identity in the
interior of the ordered time domain.  Formula \eqref{co:raw-coefficient}
will also specify its smooth extension to the faces.

\begin{lemma}
\label{co:regularity}
The forms in \eqref{co:forms} are well defined, independent of the period
cut, and smooth up to every simplex face by one-sided extension.  They
are invariant under simultaneous translation of all times, under the
centralizer action, and under twisted rotation.  Their tangential
restrictions to faces equal the forms for the shortened tuples, using the
same parameters and the same fixed integer $n$.
\end{lemma}

\begin{proof}
We first treat $j>0$.  On a fixed tuple the smoothed probabilities have
the explicit expression
\begin{equation}
 P_j(u)=\sum_{a\in\Z}\sum_{i=0}^k
 \left(\int_{t_i+a}^{t_{i+1}+a}\eta_j(u-v)\,dv\right)\delta_{g^ah_i}.
 \label{co:smooth-intervals}
\end{equation}
For $u$ in a bounded interval and times in a neighborhood of any point of
the closed ordered domain, only a fixed finite set of summands occurs.
Each coefficient is smooth in $u$ and the times, including when the
endpoints coincide.  The channel maps are smooth on the finite
probability simplices, with smooth extensions at their faces, so $Y_j$
has the same regularity.  Its possible labels on one period lie in a
fixed finite set locally in the time variables.

When $U_j$ has compact support, the two integrals in
\eqref{co:forms} can therefore be differentiated under fixed integration
domains.  The final full-sign kernel in case (i) requires an additional
observation.  Write $v=w+a$ with $w$ in one period.  If $y,z$ are possible
labels in that period, then
\[
 K_{b_*}(y,g^az)\ne0\quad\Longleftrightarrow\quad y=g^az.
\]
Since $g$ has infinite order, this equation has at most one solution
$a\in\Z$ for each pair $(y,z)$.  There are only finitely many such pairs
locally.  The full-sign integral is therefore locally a finite sum of
smooth integrals over bounded domains, even though the possible matching
integers need not be uniformly bounded as the tuple varies.

To check the period cut, regard
\[
 F_j(u)=\int_{\R}U_j(v-u)K_j(dY_j(u),Y_j(v))\,dv
\]
as a one-form in the time parameters.  The identities
$Y_j(u+1)=gY_j(u)$ and $dY_j(u+1)=g\,dY_j(u)$, followed by the change of
variable $v\mapsto v+1$, give $F_j(u+1)=F_j(u)$.  Hence its integral over
one period is independent of the cut.  This also shows why a moving cut
does not introduce an extra exterior-derivative term: if the cut is
$c=c(t)$, the endpoint contribution is
\begin{equation}
 dc\wedge\bigl(F_j(c+1)-F_j(c)\bigr)=0.
 \label{co:cut-cancellation}
\end{equation}

For the raw stage, work first away from the faces and choose a fixed
period cut away from its finitely many jumps.  Evaluating the impulse
measure gives \eqref{co:raw-coefficient}; moving an impulse to another
period leaves this coefficient unchanged, by equivariance and a change
of variable.  To prove one-sided smoothness explicitly, let
\[
 Q(s)=\int_0^sU_0(w)\,dw
 =|s|-2\int_0^s\int_0^w\rho_0(a)\,da\,dw.
\]
Then $(B_0)_i$ is a locally finite sum of terms
\begin{equation}
 -\frac12K_0(J_i,\delta_{g^ah_q})
 \bigl(Q(t_{q+1}+a-t_i)-Q(t_q+a-t_i)\bigr),
 \quad 0\leq q\leq k.
 \label{co:raw-primitive}
\end{equation}
Although $Q$ is constant rather than zero outside the lag support, each
displayed endpoint difference vanishes when its interval misses that
support.  Since all boundaries lie in one period and $2\delta_1<1$, only
the shifts $a=-1,0,1$ can contribute, with the same finite choice valid
in a sufficiently small neighborhood of any point of the closed domain.
On the entire closed ordered domain, each affine expression
$t_q+a-t_i$ has a fixed weak sign: the order of the boundaries and all
their period translates is prescribed.  The same is true for
$t_{q+1}+a-t_i$.  Thus every absolute value in
\eqref{co:raw-primitive} equals a fixed signed affine function on that
domain.  The remaining part of $Q$ is smooth on the whole line.
Consequently the coefficient has a smooth extension from the prescribed
side at each face, including intersections of faces.  Compact support of
$U_0$ allows the finite set of summands to be fixed in a neighborhood.

It remains to identify the restrictions.  On the face $t_i=t_{i+1}$,
where $0\leq i<k$, the interval labelled $h_i$ disappears and the
adjoining impulses combine tangentially as
\[
 J_i\,dt_i+J_{i+1}\,dt_{i+1}
 =\bigl(\delta_{h_{i-1}}-\delta_{h_{i+1}}\bigr)dt_i.
\]
This is precisely the impulse for the shortened tuple.  On the closing
face $t_k=t_0+1$, combine the impulses at that common lift:
\[
 J_k+gJ_0
 =\delta_{h_{k-1}}-\delta_{gh_0}.
\]
Their time differentials agree tangentially, and the displayed sum is
the closing impulse after deleting $h_k$.  These identities also handle
corners by successive merging.  In the smoothed stages, the underlying
step function itself is unchanged by deleting a zero-length interval;
\eqref{co:smooth-intervals} then gives equality of the fields and their
tangential differentials.  Hence every $B_j$ respects the faces.

Finally, translating all times by $s$ replaces $Y_j(u)$ by $Y_j(u-s)$.
Changing both integration variables by $s$ and using cut independence
proves translation invariance.  The centralizer action preserves the
twist and all pairings.  Twisted rotation describes the very same step
function with a different first boundary, so it too preserves the
forms.  These arguments also establish cut independence of the raw form
on the faces, by continuity.
\end{proof}

\subsection{Contraction and the positive normalizer}

The forms are now compatible on all tuple simplices.  To normalize them
to a connection, we must show that their sum has positive contraction
with the translation vector $V$.  The choice of the lag derivative in
\eqref{co:lag-derivative} converts this contraction into a difference
between a pairing at one time and its average at the next scale.

\begin{lemma}
\label{co:contraction-lemma}
For every stage $0\leq j\leq b_*$,
\begin{equation}
 C_j:=B_j(V)
 =\int_{\rm per}\left(
 K_j(Y_j(u),Y_j(u))
 -\int_{\R}K_j(Y_j(u),Y_j(u+s))\,d\rho_j(s)\right)du.
 \label{co:contraction}
\end{equation}
For $j<b_*$, the second integrated pairing is also
\begin{equation}
 \int_{\rm per}
 K_j\bigl((\mu_{j+1}*Y_j)(u),(\mu_{j+1}*Y_j)(u)\bigr)\,du.
 \label{co:averaged-pairing}
\end{equation}
\end{lemma}

\begin{proof}
Translation of the times translates the field in the opposite direction,
so $dY_j(V)=-\partial_uY_j$.  Define the scalar function
\[
 H_j(u)=\int_{\R}U_j(v-u)K_j(Y_j(u),Y_j(v))\,dv.
\]
As in the proof of Lemma~\ref{co:regularity}, it is periodic.  For a
smooth stage its distributional derivative is
\[
 H_j'(u)
 =-\int U_j'(v-u)K_j(Y_j(u),Y_j(v))\,dv
  +\int U_j(v-u)K_j(\partial_uY_j(u),Y_j(v))\,dv.
\]
Its derivative integrates to zero over a period.  Substitution into
\eqref{co:forms}, followed by \eqref{co:lag-derivative}, gives
\eqref{co:contraction}, with the factor $1/2$ cancelling the factor two
in $U_j'$.

The same computation is valid for $Y_0$ with its jumps.  Indeed, for
each label $y$, the function
\[
 u\longmapsto\int U_0(v-u)K_0(\delta_y,P_0(v))\,dv
\]
is continuous and has bounded weak derivative obtained by replacing
$U_0$ by $-U_0'$.  Integration by parts against the jump measure is
therefore valid.  A cut away from the jumps makes the boundary
cancellation ordinary periodicity; the result extends to the faces by
Lemma~\ref{co:regularity}.

For \eqref{co:averaged-pairing}, expand both convolutions.  If their lags
are $a,b$, simultaneous translation in the period integral changes
$K_j(Y_j(u-a),Y_j(u-b))$ into
$K_j(Y_j(u),Y_j(u+a-b))$.  This is legitimate because the scalar pairing
is invariant under simultaneous integer shifts of its two time
arguments.  The distribution of $a-b$ is
$\mu_{j+1}*\mu_{j+1}=\rho_j$, since $\mu_{j+1}$ is even.  This gives the
asserted identity.
\end{proof}

Write
\[
 \Phi_j=\int_{\rm per}\phi_{r_j}(P_j(u))\,du
 \qquad(1\leq j\leq b_*).
\]
These are the concentration quantities that will telescope.

\begin{lemma}
\label{co:positivity}
Under \eqref{co:twist-assumption}, the scalar function
\begin{equation}
 C_{\rm nor}:=B(V)=\sum_{j=0}^{b_*-1}C_j+WC_{b_*}
 \label{co:normalizer}
\end{equation}
satisfies $C_{\rm nor}\geq c_0$, where $c_0=1/10$ is independent of all
tuples, $n$, $R$, and $g$.
\end{lemma}

\begin{proof}
Consider first $1\leq j<b_*$.  The concentration-loss calculation
\eqref{ch:averaging-drop} takes the following form for our time averages.
In this paragraph abbreviate
$A_j(P_j(u))$ to $A(u)$ and the prediction
$\operatorname{pred}_j(P_j(u),y)$ to $p(u,y)$.  Put
$\overline A=\mu_{j+1}*A$ and
$\overline p(\cdot,y)=\mu_{j+1}*p(\cdot,y)$.  Equations
\eqref{co:channel-pairing} and \eqref{co:averaged-pairing} give
\[
 C_j=\int_{\rm per}\sum_y A(u)_y p(u,y)\,du
      -\int_{\rm per}\sum_y\overline A(u)_y\overline p(u,y)\,du.
\]
The total mass of $A(u)$ is a periodic scalar function.  Averaging
preserves its integral, so we may subtract $a_*$ from both predictions:
\begin{equation}
 C_j=\int_{\rm per}\sum_y A(u)_y(p(u,y)-a_*)\,du
      -\int_{\rm per}\sum_y\overline A(u)_y
                         (\overline p(u,y)-a_*)\,du.
 \label{co:subtract-threshold}
\end{equation}
This cancellation uses the integrated mass; it does not require either
channel to have mass one at each time.

The first and third inequalities of Proposition~\ref{ch:summary} imply
\[
 \sum_y A(u)_y(p(u,y)-a_*)
 \geq\phi_{r_j}(P_j(u))-20\epsilon_j.
\]
For the other term, the prediction upper bound, linearity of
$E_{r_{j+1}}$ in its probability argument, and
$P_{j+1}=\mu_{j+1}*P_j$ imply
\[
 \overline p(u,y)-a_*
 \leq E_{r_{j+1}}(P_{j+1}(u),y)-a_*+10\epsilon_j.
\]
Here the radius is $r_{j+1}=6r_j$ for the early version and
$r_{j+1}=r_j$ for the later version, exactly as required by the two
channel bounds.  Since $\overline A(u)$ is a subprobability distribution,
pairing this inequality with it yields an upper bound
$\phi_{r_{j+1}}(P_{j+1}(u))+10\epsilon_j$.  Consequently
\begin{equation}
 C_j\geq\Phi_j-\Phi_{j+1}-30\epsilon_j
 \qquad(1\leq j<b_*).
 \label{co:telescope}
\end{equation}

The raw stage supplies the initial positive quantity.  Its same-time
pairing is one, while \eqref{co:averaged-pairing} gives
$K_0(P_1(u),P_1(u))$ for the averaged pairing.  For every probability
distribution $P$,
\[
 K_0(P,P)=\sum_yP_yE_R(P,y)\leq q_R(P)
 \leq a_*+\phi_R(P).
\]
Therefore
\begin{equation}
 C_0\geq1-a_*-\Phi_1=\frac15-\Phi_1.
 \label{co:initial-contraction}
\end{equation}

At the final stage the negative term in \eqref{co:contraction} vanishes.
This is immediate from $\rho_{b_*}=0$ in case (i).  In case (ii), let
$S=\supp D(P_{b_*}(u))$.  Lemma~\ref{ch:terminal} gives
$\operatorname{diam}S\leq4r$.  If $z=gw$ with $z,w\in S$, then
$d(w,gw)\leq4r$, contrary to \eqref{co:twist-assumption}.  Thus
$S\cap gS=S\cap g^{-1}S=\varnothing$, and twist equivariance makes both
lag pairings at $s=\pm1$ zero.  In either case,
\[
 C_{b_*}=\int_{\rm per}\sum_zD(P_{b_*}(u))_z^2\,du,
 \qquad WC_{b_*}\geq\Phi_{b_*}
\]
by \eqref{co:terminal-multiplier}.  Adding this inequality to
\eqref{co:telescope} and \eqref{co:initial-contraction} leaves
\[
 C_{\rm nor}\geq\frac15-30\sum_{j=1}^{b_*-1}\epsilon_j
 \geq\frac1{10},
\]
as claimed.
\end{proof}

\subsection{Differential bounds and the resulting connection}

Positivity now allows the normalization.  Before applying the trace
identity, we also need growth bounds that justify summation over all
tuples.  More precise bounds on the individual channel differentials
will provide the input for the curvature estimates of the next section.

We now specialize the parameters of the preceding construction.  At the
shortest time scales we use the early channels, whose derivative bound
is independent of the spatial radius.  The later channels keep that
radius fixed; this also fixes the terminal multiplier $W$ while the
time scales continue to increase.  Their longer time scales will absorb
the polynomial factor $r^{10}$ in the later derivative bound.

Choose once and for all an integer $J_*$ with
$(3/2)^{J_*-1}>200$.  For sufficiently large $R$ retain the spatial
radii, accuracies, and multiplier from
\eqref{co:spatial-radii}, \eqref{co:error-budget}, and
\eqref{co:accuracy}, and set
\begin{equation}
 \beta_1=R,\qquad \beta_{j+1}=\beta_j^{3/2}.
 \label{co:ladder}
\end{equation}
Choose $b_*$ to be the smallest integer such that
\begin{equation}
 b_*\geq J_*,\qquad \beta_{b_*}\geq W^2R^{100}.
 \label{co:last-scale}
\end{equation}
For each even integer $n>10\beta_{b_*}$, use
\begin{equation}
 \delta_j=\beta_j/n\qquad(1\leq j\leq b_*).
 \label{co:time-scales}
\end{equation}
The ratio $\beta_{j+1}/\beta_j$ is at least two when $R$ is sufficiently
large, so these scales satisfy \eqref{co:scale-support}.  Thus the
construction and positivity proved above apply.  For each such $n$, the
same actual scales $\beta_j/n$ are used in every dimension $k\leq n$;
face restriction never changes the denominator to $k$.

The stage count, spatial radii, accuracies, multiplier, and numbers
$\beta_j$ now depend only on $R$ and the fixed choices.  Since
$\log\beta_j=(3/2)^{j-1}\log R$ and
$\log(W^2R^{100})=2C_*r+100\log R$, with $r$ a fixed multiple of $R$,
\begin{equation}
 b_*=O(1+\log R).
 \label{co:stage-count}
\end{equation}
The particular exponents in \eqref{co:ladder} and
\eqref{co:last-scale} will enter the curvature bounds.  They leave room
for the polynomial losses in the channel derivatives.

Choose a sufficiently large fixed constant $C$ and set
\begin{equation}
 a_j=C\epsilon_j^{-10}
 \begin{cases}
 1,&1\leq j<J_*,\\
 r^{10},&J_*\leq j\leq b_*.
 \end{cases}
 \label{co:derivative-constants}
\end{equation}
As before, $l_i=1+d(h_{i-1},h_i)$, with periodic indexing, and
$\partial_i=\partial/\partial t_i$.

\begin{proposition}
\label{co:connection}
Let $R$ be sufficiently large, choose the parameters in
\eqref{co:ladder} and \eqref{co:last-scale}, and let $n>10\beta_{b_*}$
be even, with time scales \eqref{co:time-scales}.  Suppose that $g$
satisfies \eqref{co:twist-assumption}.  The forms $B_j$ and $B$ in
\eqref{co:forms}, constructed in all dimensions $0\leq k\leq n$ with the
same parameters, have the following properties.  For every smooth stage
$1\leq j\leq b_*$ and every coordinate $t_i$,
\begin{equation}
 \left(\int_{\rm per}\sum_\omega
        \frac{|\partial_iY_j(u)_\omega|^2}{Y_j(u)_\omega}\,du
 \right)^{1/2}
 \leq\frac{a_jl_i}{\sqrt{\delta_j}}.
 \label{co:fisher}
\end{equation}
Terms with zero denominator are defined to be zero.  Moreover,
\begin{equation}
 |B(\partial_i)|+|\partial_iC_{\rm nor}|\leq C_Rl_i.
 \label{co:coefficients}
\end{equation}
The scalar $C_{\rm nor}$ is at least $c_0=1/10$, and
\begin{equation}
 \alpha=\frac{B}{C_{\rm nor}}
 \label{co:normalized}
\end{equation}
is a compatible connection family in the sense of
Definition~\ref{cy:connection}.  Its coefficients and those of its
exterior derivative are bounded by polynomials in the lengths, uniformly
over the tuples.  In particular, the family may be used in
Proposition~\ref{cy:trace-identity}.
\end{proposition}

\begin{proof}
Differentiating \eqref{co:smooth-intervals} at a boundary gives the smeared
impulse formula
\begin{equation}
 \partial_iP_j(u)
 =\sum_{a\in\Z}\eta_j(u-t_i-a)g^aJ_i.
 \label{co:smeared-impulse}
\end{equation}
For each $u$, at most one summand can be nonzero, because
$4\delta_j<1$.  The derivative is supported within $2\delta_j$ of
$t_i$ modulo periods.  The distance between the two labels in $g^aJ_i$
is $d(h_{i-1},h_i)$, independently of $a$.  The derivative estimate of
Proposition~\ref{ch:summary}, or Lemma~\ref{ch:terminal} at the last
stage, therefore bounds the square root of the pointwise sum in
\eqref{co:fisher} by
$C\epsilon_j^{-10}l_i\eta_j(u-t_i-a)$ for $j<J_*$, with an additional
factor $r^{10}$ for $j\geq J_*$, where $a$ is the contributing integer.
Since
\[
 \int_{\rm per}\sum_a\eta_j(u-t_i-a)^2\,du
 =\int_{\R}\eta_j(s)^2\,ds
 \leq\|\eta_j\|_\infty\leq C/\delta_j,
\]
we obtain \eqref{co:fisher} after fixing the constant in
\eqref{co:derivative-constants}.  The two disjoint channels only alter
this constant.  At a zero channel weight the first derivative is zero,
as stated in Proposition~\ref{ch:summary}, so the convention in
\eqref{co:fisher} is consistent.

The support information is useful as well as the square-integral bound.
The support of $\partial_iY_j$ has measure at most $4\delta_j$ in a
period, and the field has mass at most two.  Cauchy--Schwarz in time and
labels consequently gives
\begin{equation}
 \int_{\rm per}\sum_\omega|\partial_iY_j(u)_\omega|\,du
 \leq C a_jl_i.
 \label{co:integrated-derivative}
\end{equation}
For example, the square of the factor paired with
\eqref{co:fisher} is at most
$\int_{\supp\partial_iY_j}\sum_\omega Y_j(u)_\omega\,du
\leq8\delta_j$.  This explains why the integrated derivative loses no
power of the time scale.

We now bound the coefficients of $B$.  For a compact lag kernel, the
inner integral in \eqref{co:forms} has length at most two, its kernel is
bounded, and the field has bounded mass at every time.  Thus
\eqref{co:integrated-derivative} gives
$|B_j(\partial_i)|\leq Ca_jl_i$.  At the full-sign stage there is no
bound on the length of the integration interval, but equality of labels
gives a stronger estimate.  For each fixed group label $y$,
\begin{equation}
 \int_{\R}Y_{b_*}(v)_y\,dv
 =\int_0^1\sum_{a\in\Z}Y_{b_*}(v)_{g^{-a}y}\,dv
 \leq1.
 \label{co:orbit-mass}
\end{equation}
The labels $g^{-a}y$ are distinct because $g$ has infinite order, and the
total mass at every time is at most one.  Combining
\eqref{co:orbit-mass} with \eqref{co:integrated-derivative} proves the
same coefficient bound.  In particular, the estimate does not depend
on how far apart the finitely many matching periods are.

Differentiating \eqref{co:contraction} gives
$|\partial_iC_j|\leq Ca_jl_i$ for $j>0$.  Indeed, each differentiation
puts a derivative on one of the two factors.  The other factor has
bounded total mass, the measure $\rho_j$ has mass at most one, and
\eqref{co:integrated-derivative} is unchanged by a time shift.  The last
observation follows because the sum of absolute derivatives is a
periodic scalar function.

For the raw stage, \eqref{co:raw-coefficient} bounds
$|(B_0)_i|\leq C$, since $\|J_i\|_1\leq2$ and the lag support has length
at most $4\delta_1<1$.  Its contraction satisfies the exact identity
\[
 C_0=1-\int_{\rm per}K_0(P_1(u),P_1(u))\,du.
\]
Formula \eqref{co:smeared-impulse} for $P_1$, before passing through any
channel map, gives
$\int_{\rm per}\|\partial_iP_1(u)\|_1\,du\leq2$.
Hence $|\partial_iC_0|\leq4$.  Summing these estimates with the
coefficients $1,\ldots,1,W$ proves \eqref{co:coefficients}, because
\[
 1+\sum_{j=1}^{b_*-1}a_j+Wa_{b_*}
\]
depends only on $R$ and the fixed choices.

For completeness, polynomial control of exterior derivatives can be
obtained here without using the finer estimates of the next section.
For a smooth stage, differentiation at a fixed period cut gives
\begin{equation}
 dB_j=\frac12\int_{\rm per}du\int_{\R}dv\,
          U_j(v-u)K_j(dY_j(u)\wedge dY_j(v)).
 \label{co:exterior-pairing}
\end{equation}
There is no term from the cut by \eqref{co:cut-cancellation}.  For a
compact lag kernel, reduction of both time variables to one period
introduces at most three relevant period shifts.  Boundedness of the
kernel and \eqref{co:integrated-derivative} therefore imply
\[
 |dB_j(\partial_i,\partial_{i'})|
 \leq Ca_j^2l_il_{i'}.
\]
For the full-sign equality kernel, at most one shift contributes for
each pair of labels, and the same estimate follows.  In the raw stage,
differentiate \eqref{co:raw-coefficient} with respect to a distinct
boundary.  The resulting terms have factors
$U_0(t_{i'}+a-t_i)K_0(J_i,g^aJ_{i'})$; boundedness of $U_0,K_0$ and the
jump masses gives a uniform bound.  Differentiation of a coefficient in
its own time variable contributes no exterior two-form.  These
calculations hold in the interior and extend to the closed simplex by
Lemma~\ref{co:regularity}.  Summation over stages gives
\begin{equation}
 |dB(\partial_i,\partial_{i'})|\leq C_Rl_il_{i'}.
 \label{co:coarse-curvature}
\end{equation}

Finally Lemma~\ref{co:positivity} permits division by $C_{\rm nor}$.
The scalar normalizer inherits all symmetries and face restrictions of
$B$, so \eqref{co:normalized} does also, and $\alpha(V)=1$.
Equations \eqref{co:coefficients}, \eqref{co:coarse-curvature}, and
\[
 d\alpha=C_{\rm nor}^{-1}dB
          -C_{\rm nor}^{-2}dC_{\rm nor}\wedge B
\]
give the required polynomial coefficient bounds.  The resulting family
is therefore a compatible connection family.  Its curvature is basic, as
also follows directly from translation invariance and $\alpha(V)=1$.
\end{proof}

The construction separates two estimates that the trace argument will
use differently.  The bounds in Proposition~\ref{co:connection} justify
the weighted simplex identities, while the scale-dependent bound
\eqref{co:fisher} retains the information needed to make high-dimensional
curvature Pfaffians small.  We turn to those sharper estimates next.
